\chapter{引子}
\label{chap:introduction}

理论计算机科学经常把两个层面分开讨论：函数（function）描述要计算
\emph{什么}，是问题的规格；程序（program）描述\emph{如何}计算，是规格
的一种实现。本讲先用二进制串形式化“描述”，再比较函数与程序在数量上的
差异。

全文约定
\[
  \Nat=\{0,1,2,\ldots\},\qquad \Sigma=\bits.
\]

\section{二进制串与编码}

\begin{definition}[二进制串]
  \label{def:binary-string}
  对任意 $n\in\Nat$，定义
  \[
    \Sigma^n=\{(a_1,\ldots,a_n):a_i\in\Sigma\}.
  \]
  特别地，$\Sigma^0=\{\eps\}$，其中 $\eps$ 表示空串。所有有限二进制串的
  集合记为
  \[
    \Sigma^*=\bigcup_{n\geq 0}\Sigma^n,
  \]
  非空二进制串的集合记为 $\Sigma^+=\Sigma^*\setminus\{\eps\}$。若
  $x\in\Sigma^n$，则称 $x$ 的长度为 $\card{x}=n$。

  对 $x,y\in\Sigma^*$，把二者的串联记为 $xy$。若存在 $z\in\Sigma^*$
  使 $y=xz$，则称 $x$ 是 $y$ 的\emph{前缀}；若还能取到
  $z\neq\eps$，则称 $x$ 是 $y$ 的\emph{真前缀}。
\end{definition}

\begin{definition}[编码]
  \label{def:encoding}
  集合 $A$ 上的一个（二进制）\emph{编码}是单射
  \[
    E:A\longrightarrow\Sigma^*.
  \]
  单射条件保证每个码字至多对应 $A$ 中的一个对象，因而解码时不会混淆。
\end{definition}

\begin{example}[直接串联可能产生歧义]
  \label{ex:pair-encoding}
  令 $b(n)$ 为自然数 $n$ 的标准二进制表示，并尝试用
  \[
    P(a,b)=b(a)b(b)
  \]
  编码 $\Nat\times\Nat$。这个映射不是单射。例如
  \[
    P(1,6)=1\,110=1110=11\,10=P(3,2).
  \]
  问题在于直接串联后无法确定两个分量的分界位置。
\end{example}

\begin{definition}[无前缀编码]
  \label{def:prefix-free}
  编码 $E:A\to\Sigma^+$ 称为\emph{无前缀的}（prefix-free），如果对任意
  不同的 $x,y\in A$，$E(x)$ 都不是 $E(y)$ 的前缀。
\end{definition}

这里要求码字非空不是多余的：如果唯一的码字是 $\eps$，那么串联任意多个
码字都仍然得到 $\eps$，不可能唯一解码。

记 $A^*$ 为由 $A$ 中元素组成的所有有限序列的集合，其空序列仍记为
$\eps$。对映射 $E:A\to\Sigma^*$，定义其逐项串联扩张
\begin{equation}
  \label{eq:concatenation-extension}
  \overline E(a_1,\ldots,a_n)=
  \begin{cases}
    E(a_1)\cdots E(a_n), & n>0,\\
    \eps, & n=0.
  \end{cases}
\end{equation}

\begin{theorem}[无前缀编码的唯一可解码性]
  \label{thm:prefix-extension-injective}
  若 $E:A\to\Sigma^+$ 是无前缀编码，则由
  \cref{eq:concatenation-extension} 定义的
  $\overline E:A^*\to\Sigma^*$ 是单射。
\end{theorem}

\begin{proof}
  假设
  \[
    \overline E(a_1,\ldots,a_p)
    =\overline E(b_1,\ldots,b_q).
  \]
  因为每个码字非空，等式一侧为空串当且仅当对应序列为空；故只需考虑
  $p,q>0$ 的情形。两个相等的串分别以 $E(a_1)$ 和 $E(b_1)$ 开头，因此
  $E(a_1)$ 与 $E(b_1)$ 中较短者必为较长者的前缀。由无前缀性可知
  $E(a_1)=E(b_1)$，再由 $E$ 的单射性得到 $a_1=b_1$。

  从等式两侧消去相同的首个码字并重复上述论证，最终得到 $p=q$ 且
  $a_i=b_i$（$1\leq i\leq p$）。因此 $\overline E$ 是单射。
\end{proof}

\begin{theorem}[编码的无前缀化]
  \label{thm:prefix-free-construction}
  若存在编码 $E:A\to\Sigma^*$，则存在无前缀编码
  $E_{\mathrm{pf}}:A\to\Sigma^+$，并且对每个 $a\in A$ 都有
  \[
    \card{E_{\mathrm{pf}}(a)}=2\card{E(a)}+2.
  \]
\end{theorem}

\begin{proof}
  先定义 $h(0)=00$、$h(1)=11$，并令 $h$ 在串上逐位作用。取
  \[
    E_{\mathrm{pf}}(a)=h(E(a))01.
  \]
  其长度显然为 $2\card{E(a)}+2$。

  把 $E_{\mathrm{pf}}(a)$ 从左到右按两个比特一组切分：最后一组是
  $01$，此前每一组只能是 $00$ 或 $11$。假设
  $E_{\mathrm{pf}}(a)$ 是 $E_{\mathrm{pf}}(b)$ 的前缀。由于两个串的长度
  都是偶数，短串结尾的 $01$ 与长串中的某个完整二比特组对齐；它不可能
  等于长串正文中的 $00$ 或 $11$，所以只能与长串末尾的 $01$ 对齐。
  因而两个码字等长且相等。逐组解码可得 $E(a)=E(b)$，再由 $E$ 的单射性
  得 $a=b$。所以不存在一个码字是另一个不同码字的前缀。
\end{proof}

\begin{remark}
  原稿进一步询问能否把长度开销改进到
  $\card{E(a)}+O(\log \card{E(a)})$。答案是肯定的。令
  $n=\card{E(a)}$，先用自定界格式
  \[
    1^{\ell}0\,\operatorname{bin}(n+1),
    \qquad \ell=\card{\operatorname{bin}(n+1)},
  \]
  写出长度，再接上 $E(a)$。解码器先由连续的 $1$ 得到 $\ell$，继而读出
  $n+1$，最后恰好读取 $n$ 位正文。所得码长为
  \[
    n+2\lfloor\log_2(n+1)\rfloor+3=n+O(\log(n+1)),
  \]
  且构造仍然无前缀。
\end{remark}

\begin{corollary}
  \label{cor:sequence-encoding}
  若存在编码 $E:A\to\Sigma^*$，则存在编码
  $F:A^*\to\Sigma^*$。
\end{corollary}

\begin{proof}
  先用\cref{thm:prefix-free-construction}把 $E$ 无前缀化，再用
  \cref{thm:prefix-extension-injective}逐项串联即可。
\end{proof}

\section{可数集合}

\begin{definition}[至多可数]
  \label{def:countable}
  若存在单射 $f:A\to\Nat$，则称集合 $A$ 是\emph{至多可数的}。这等价于
  $A$ 有限或 $A$ 与 $\Nat$ 之间存在双射。若 $A\neq\varnothing$，它也
  等价于存在满射 $g:\Nat\to A$。
\end{definition}

\begin{remark}
  非空条件只影响最后一种表述：不存在从 $\Nat$ 到空集的函数，但空集显然
  是至多可数的。部分文献把“可数”专指与 $\Nat$ 等势，把本文的“至多
  可数”称作“可数”；使用时应先确认约定。
\end{remark}

\begin{lemma}
  \label{lem:binary-strings-countable}
  $\Sigma^*$ 是可数无限集。
\end{lemma}

\begin{proof}
  对长度为 $n$ 的二进制串 $x$，令 $\operatorname{val}_2(x)$ 表示允许前导
  零时的二进制数值，并定义
  \[
    \nu(x)=2^n-1+\operatorname{val}_2(x).
  \]
  长度为 $n$ 的串恰好映到区间
  $\{2^n-1,\ldots,2^{n+1}-2\}$；这些区间互不相交且覆盖 $\Nat$。
  因此 $\nu:\Sigma^*\to\Nat$ 是双射。
\end{proof}

\begin{theorem}[可编码性与可数性]
  \label{thm:encodable-iff-countable}
  集合 $A$ 至多可数，当且仅当存在编码 $E:A\to\Sigma^*$。
\end{theorem}

\begin{proof}
  若 $E:A\to\Sigma^*$ 是编码，则把它与
  \cref{lem:binary-strings-countable}中的单射
  $\nu:\Sigma^*\to\Nat$ 复合，便得到 $A\to\Nat$ 的单射，所以 $A$
  至多可数。

  反之，若 $f:A\to\Nat$ 是单射，则把 $f$ 与映射
  $n\mapsto 1^n0$ 复合，得到 $A\to\Sigma^*$ 的单射，也就是一个编码。
\end{proof}

\section{问题比程序多}

\begin{definition}[问题]
  \label{def:problem}
  在本讲义中，一个（计算）\emph{问题}是函数
  \[
    f:\Sigma^*\longrightarrow\Sigma^*.
  \]
  所有问题构成的集合记为
  \[
    \mathcal{P}=(\Sigma^*)^{\Sigma^*}.
  \]
\end{definition}

\begin{theorem}
  \label{thm:problems-uncountable}
  问题集合 $\mathcal{P}$ 不可数。
\end{theorem}

\begin{proof}[证明（Cantor 对角化）]
  由\cref{lem:binary-strings-countable}，可将所有输入依次写成
  $s_0,s_1,s_2,\ldots$。假设所有问题也能被序列
  $f_0,f_1,f_2,\ldots$ 穷尽。构造新问题 $g:\Sigma^*\to\Sigma^*$：
  \[
    g(s_i)=
    \begin{cases}
      0, & f_i(s_i)\neq 0,\\
      1, & f_i(s_i)=0.
    \end{cases}
  \]
  对每个 $i\in\Nat$ 都有 $g(s_i)\neq f_i(s_i)$，所以 $g\neq f_i$。
  这与该序列穷尽所有问题的假设矛盾。因此 $\mathcal{P}$ 不可数。
\end{proof}

\begin{corollary}
  \label{cor:uncomputable-problems}
  如果程序集合存在一个有限二进制编码，并且每个程序至多实现一个问题，那么
  存在没有任何程序能够实现的问题。
\end{corollary}

\begin{proof}
  由\cref{thm:encodable-iff-countable}，所有程序构成至多可数集，因而程序
  能实现的问题也至多可数；但\cref{thm:problems-uncountable}说明全部问题
  的集合不可数。
\end{proof}
